% !TeX program = lualatex
% ================================================================
% Polish Parallel Text Version of PercentagesQuickQuestions
%
% Generated by the server using the following settings:
%
%   language            Polish (pol)
%   text direction      left to right
%   font                Noto Serif
%   fallback font       Noto Serif
%   built from          latex/quickQuestions/percentagesQuickQuestions.tex
%   vocabulary key      latex/quickQuestions/percentagesQuickQuestions/L2/pol/percentagesQuickQuestions_polishKey.tex
%   tier 3 vocabulary   green   RGB 0 128 0
%   English text        black   RGB 0 0 0
%   Polish text         purple  RGB 112 48 160
%   layout macros       version 3
%
% Please don't edit any information in this comment block - the server
% compares it against the current settings and rebuilds this file when
% they differ, so changes here would be overwritten. However, feel free
% to make updates to the LaTeX below if you want to edit it.
% ================================================================

\documentclass{beamer}
% ================================================================
% Copied from the English sheet this was translated from, so that
% anything it sets up for itself still works here
% ================================================================
\usepackage{amsmath}
\usepackage{amssymb}
\usepackage{tikz}
\usetikzlibrary{decorations.pathreplacing}

% ================================================================
% Quick Questions deck: percentages.
%
% Every question gets two slides: the question on its own for the
% mini whiteboards, then the same question again with the solution
% underneath in red. Within each topic the ten questions get
% gradually harder.
% ================================================================

\setbeamertemplate{navigation symbols}{}
\setbeamertemplate{footline}{}

% Topic divider.
\newcommand{\qtopic}[1]{%
  \section{#1}%
  \begin{frame}[plain]
    \vfill
    \begin{center}
      {\usebeamercolor[fg]{structure}\Huge\bfseries #1}
    \end{center}
    \vfill
  \end{frame}%
}

% \qq[<question size>][<solution size>]{<question>}{<solution>}
% The two optional sizes drop the type for the wordier questions and
% the longer solutions, so nothing runs off the slide.
\NewDocumentCommand{\qq}{O{\Huge}O{\LARGE}+m+m}{%
  \begin{frame}
    \vfill
    \begin{center}
      \begin{minipage}{0.92\textwidth}
        \centering #1 #3
      \end{minipage}
    \end{center}
    \vfill
  \end{frame}%
  \begin{frame}
    \vfill
    \begin{center}
      \begin{minipage}{0.92\textwidth}
        \centering #1 #3
        \par\vspace{0.9em}
        {\color{red}#2 #4}
      \end{minipage}
    \end{center}
    \vfill
  \end{frame}%
}

% Stepped working, written in the form
%   80% is 64 / 10% is 8 / 100% is 80
% Rows are separated with \ and the two columns with &.
\newcommand{\steps}[1]{%
  \begin{tabular}{@{}r@{\ is\ }l@{}}#1\end{tabular}%
}

% \barmodel{<cells>}{<cells spanned by the lower brace>}{<lower label>}{<upper label>}
% The whole bar is 9cm wide and represents the upper label; the lower
% brace picks out the part of it named by the lower label.
\newcommand{\barmodel}[4]{%
  \begin{tikzpicture}[x=1cm,y=1cm]
    \pgfmathsetmacro{\cw}{9/#1}
    \fill[red!12] (0,0) rectangle ({#2*\cw},0.7);
    \foreach \i in {1,...,#1}
      {\draw[red!60!black] ({(\i-1)*\cw},0) rectangle ({\i*\cw},0.7);}
    \draw[red!60!black,thick] (0,0) rectangle (9,0.7);
    \draw[red,decorate,decoration={brace,amplitude=4pt}] (0,0.9) -- (9,0.9)
      node[midway,above=4pt,red,font=\normalsize]{#4};
    \draw[red,decorate,decoration={brace,mirror,amplitude=4pt}] (0,-0.2) -- ({#2*\cw},-0.2)
      node[midway,below=4pt,red,font=\normalsize]{#3};
  \end{tikzpicture}%
}

% Contents-slide bullets, colour-coded by difficulty band:
% green for the first two topics, orange for the next two,
% red for the last three.
\setbeamertemplate{section in toc}{%
  \ifnum\inserttocsectionnumber<3
    \textcolor{green!55!black}{$\bullet$}%
  \else\ifnum\inserttocsectionnumber<5
    \textcolor{orange}{$\bullet$}%
  \else
    \textcolor{red}{$\bullet$}%
  \fi\fi
  ~\inserttocsection%
}

\title{Percentages: Quick Questions}
\subtitle{Mini whiteboard questions, with solutions}
\author{}
\date{}

% Characters this language's font has not got are borrowed from another,
% rather than being silently dropped - see docs/translated-sheet-latex.md
\directlua{%
  if luaotfload and luaotfload.add_fallback then
    luaotfload.add_fallback("eallatinfallback",
      { "Noto Serif:mode=harf;" })
  else
    tex.error("this luaotfload is too old to register a fallback font")
  end
}
\usepackage[english]{babel}
\babelprovide[import]{polish}
\babelfont[polish]{rm}[Renderer=Harfbuzz, BoldFont={Noto Serif}, ItalicFont={Noto Serif}, BoldItalicFont={Noto Serif}, RawFeature={fallback=eallatinfallback}]{Noto Serif}
\babelfont[polish]{sf}[Renderer=Harfbuzz, BoldFont={Noto Serif}, ItalicFont={Noto Serif}, BoldItalicFont={Noto Serif}, RawFeature={fallback=eallatinfallback}]{Noto Serif}
\newcommand{\ealtext}[1]{\foreignlanguage{polish}{#1}}
\newcommand{\ealtextblock}[1]{{\raggedright\foreignlanguage{polish}{#1}\par}}
% ================================================================
% EAL layout helpers
% ================================================================
\usepackage{xcolor}

\definecolor{ealkeycolour}{RGB}{0,128,0}
\definecolor{ealenglish}{RGB}{0,0,0}
\definecolor{ealtwocolour}{RGB}{112,48,160}

\newsavebox{\ealboxen}
\newsavebox{\ealboxtr}
\newlength{\ealglosswd}
\newlength{\ealglossraise}
\setlength{\ealglossraise}{1.15em}

% a tier 3 word, in English and in translation alike
\newcommand{\ealkey}[1]{\textcolor{ealkeycolour}{#1}}
\newcommand{\ealkeytr}[1]{\textcolor{ealkeycolour}{\ealtext{#1}}}

% One translated word above one English word. Built from kernel
% primitives, never a tabular, which would break wherever the sheet
% loads array, booktabs or tabularx. The box is as wide as the wider of
% the two, so a long translation pushes its neighbours aside instead of
% overlapping them, and it claims its full height so a table row or a
% TikZ node makes room for it.
\newcommand{\ealgloss}[2]{%
  \sbox{\ealboxen}{\ealkey{#1}}%
  \sbox{\ealboxtr}{\small\textcolor{ealtwocolour}{\ealtext{#2}}}%
  \setlength{\ealglosswd}{\wd\ealboxen}%
  \ifdim\wd\ealboxtr>\ealglosswd\setlength{\ealglosswd}{\wd\ealboxtr}\fi
  \leavevmode
  \makebox[\ealglosswd][c]{%
    \makebox[0pt][c]{\usebox{\ealboxen}}%
    \makebox[0pt][c]{\raisebox{\ealglossraise}{\usebox{\ealboxtr}}}%
  }%
}

% opens up line spacing around a block of glosses, so the raised
% translations sit clear of the line above
\newenvironment{ealglossed}{\par\linespread{2.1}\selectfont\ignorespaces}{\par}

% a whole translated sentence above its English counterpart. block level,
% so it does not belong inside a tabular cell
\newcommand{\ealpara}[2]{%
  \par\addvspace{0.5\baselineskip}%
  {\color{ealtwocolour}\ealtextblock{#1}}%
  \penalty200
  {\color{ealenglish}#2\par}%
  \addvspace{0.5\baselineskip}%
}

\begin{document}

\frame{\titlepage}

\begin{frame}{Contents}
  \ealpara{Spis treści}{Contents}
  \small
  \tableofcontents
\end{frame}

%================================================================
\section{Converting between decimals and percentages}
\begin{frame}[plain]
  \vfill
  \begin{center}
    {\usebeamercolor[fg]{structure}\Huge\bfseries
    \ealpara{Konwersja między \ealkeytr{ułamkami dziesiętnymi} a \ealkeytr{procentami}}{Converting between \ealkey{decimals} and \ealkey{percentages}}
    }
  \end{center}
  \vfill
\end{frame}
%================================================================

\qq[\LARGE]{\ealpara{Zapisz $0.5$ jako \ealkeytr{procent}}{Write $0.5$ as a \ealkey{percentage}}}{$0.5\times 100=50\%$}
\qq[\LARGE]{\ealpara{Zapisz $0.25$ jako \ealkeytr{procent}}{Write $0.25$ as a \ealkey{percentage}}}{$0.25\times 100=25\%$}
\qq[\LARGE]{\ealpara{Zapisz $30\%$ jako \ealkeytr{ułamek dziesiętny}}{Write $30\%$ as a \ealkey{decimal}}}{$30\div 100=0.3$}
\qq[\LARGE]{\ealpara{Zapisz $0.07$ jako \ealkeytr{procent}}{Write $0.07$ as a \ealkey{percentage}}}{$0.07\times 100=7\%$}
\qq[\LARGE]{\ealpara{Zapisz $6\%$ jako \ealkeytr{ułamek dziesiętny}}{Write $6\%$ as a \ealkey{decimal}}}{$6\div 100=0.06$}
\qq[\LARGE]{\ealpara{Zapisz $0.125$ jako \ealkeytr{procent}}{Write $0.125$ as a \ealkey{percentage}}}{$0.125\times 100=12.5\%$}
\qq[\LARGE]{\ealpara{Zapisz $4.5\%$ jako \ealkeytr{ułamek dziesiętny}}{Write $4.5\%$ as a \ealkey{decimal}}}{$4.5\div 100=0.045$}
\qq[\LARGE]{\ealpara{Zapisz $1.3$ jako \ealkeytr{procent}}{Write $1.3$ as a \ealkey{percentage}}}{$1.3\times 100=130\%$}
\qq[\LARGE]{\ealpara{Zapisz $250\%$ jako \ealkeytr{ułamek dziesiętny}}{Write $250\%$ as a \ealkey{decimal}}}{$250\div 100=2.5$}
\qq[\LARGE]{\ealpara{Zapisz $0.0025$ jako \ealkeytr{procent}}{Write $0.0025$ as a \ealkey{percentage}}}{$0.0025\times 100=0.25\%$}

%================================================================
\section{Writing one number as a percentage of another}
\begin{frame}[plain]
  \vfill
  \begin{center}
    {\usebeamercolor[fg]{structure}\Huge\bfseries
    \ealpara{Zapisywanie jednej liczby jako \ealkeytr{procent} drugiej liczby}{Writing one number as a \ealkey{percentage} of another}
    }
  \end{center}
  \vfill
\end{frame}
%================================================================

\qq[\LARGE][\Large]{\ealpara{Zapisz $25$ jako \ealkeytr{procent} liczby $100$}{Write $25$ as a \ealkey{percentage} of $100$}}
  {$\dfrac{25}{100}=25\%$}
\qq[\LARGE][\Large]{\ealpara{Zapisz $7$ jako \ealkeytr{procent} liczby $10$}{Write $7$ as a \ealkey{percentage} of $10$}}
  {$\dfrac{7}{10}=\dfrac{70}{100}=70\%$}
\qq[\LARGE][\Large]{\ealpara{Zapisz $13$ jako \ealkeytr{procent} liczby $20$}{Write $13$ as a \ealkey{percentage} of $20$}}
  {$\dfrac{13}{20}=\dfrac{65}{100}=65\%$}
\qq[\LARGE][\Large]{\ealpara{Zapisz $12$ jako \ealkeytr{procent} liczby $25$}{Write $12$ as a \ealkey{percentage} of $25$}}
  {$\dfrac{12}{25}=\dfrac{48}{100}=48\%$}
\qq[\LARGE][\Large]{\ealpara{Zapisz $18$ jako \ealkeytr{procent} liczby $40$}{Write $18$ as a \ealkey{percentage} of $40$}}
  {$\dfrac{18}{40}=\dfrac{45}{100}=45\%$}
\qq[\LARGE][\Large]{\ealpara{Zapisz $21$ jako \ealkeytr{procent} liczby $50$}{Write $21$ as a \ealkey{percentage} of $50$}}
  {$\dfrac{21}{50}=\dfrac{42}{100}=42\%$}
\qq[\LARGE][\Large]{\ealpara{Zapisz $9$ jako \ealkeytr{procent} liczby $12$}{Write $9$ as a \ealkey{percentage} of $12$}}
  {$\dfrac{9}{12}=\dfrac{3}{4}=\dfrac{75}{100}=75\%$}
\qq[\LARGE][\Large]{\ealpara{Zapisz $27$ jako \ealkeytr{procent} liczby $60$}{Write $27$ as a \ealkey{percentage} of $60$}}
  {$\dfrac{27}{60}=\dfrac{9}{20}=\dfrac{45}{100}=45\%$}
\qq[\LARGE][\Large]{\ealpara{Zapisz $34$ jako \ealkeytr{procent} liczby $80$}{Write $34$ as a \ealkey{percentage} of $80$}}
  {$\dfrac{34}{80}=\dfrac{17}{40}=\dfrac{42.5}{100}=42.5\%$}
\qq[\LARGE][\Large]{\ealpara{W teście Ali zdobył $57$ z $75$.\\[0.4em] Jaki \ealkeytr{procent} zdobył Ali?}{In a test Ali scored $57$ out of $75$.\\[0.4em] What \ealkey{percentage} did Ali score?}}
  {$\dfrac{57}{75}=\dfrac{19}{25}=\dfrac{76}{100}=76\%$}

%================================================================
\section{Basic percentages without a calculator}
\begin{frame}[plain]
  \vfill
  \begin{center}
    {\usebeamercolor[fg]{structure}\Huge\bfseries
    \ealpara{Podstawowe \ealkeytr{procenty} bez kalkulatora}{Basic \ealkey{percentages} without a calculator}
    }
  \end{center}
  \vfill
\end{frame}
%================================================================

\qq{\ealpara{$50\%$ z $80$}{$50\%$ of $80$}}{$80\div 2=40$}
\qq{\ealpara{$10\%$ z $70$}{$10\%$ of $70$}}{$70\div 10=7$}
\qq{\ealpara{$10\%$ z $45$}{$10\%$ of $45$}}{$45\div 10=4.5$}
\qq{\ealpara{$1\%$ z $300$}{$1\%$ of $300$}}{$300\div 100=3$}
\qq{\ealpara{$20\%$ z $45$}{$20\%$ of $45$}}{\steps{$10\%$ & $4.5$ \\ $20\%$ & $9$}}
\qq{\ealpara{$5\%$ z $80$}{$5\%$ of $80$}}{\steps{$10\%$ & $8$ \\ $5\%$ & $4$}}
\qq{\ealpara{$30\%$ z $250$}{$30\%$ of $250$}}{\steps{$10\%$ & $25$ \\ $30\%$ & $75$}}
\qq{\ealpara{$1\%$ z $42$}{$1\%$ of $42$}}{$42\div 100=0.42$}
\qq{\ealpara{$15\%$ z $60$}{$15\%$ of $60$}}{\steps{$10\%$ & $6$ \\ $5\%$ & $3$ \\ $15\%$ & $9$}}
\qq{\ealpara{$35\%$ z $240$}{$35\%$ of $240$}}{\steps{$10\%$ & $24$ \\ $30\%$ & $72$ \\ $5\%$ & $12$ \\ $35\%$ & $84$}}

%================================================================
\section{Increasing and decreasing without a calculator}
\begin{frame}[plain]
  \vfill
  \begin{center}
    {\usebeamercolor[fg]{structure}\Huge\bfseries
    \ealpara{\ealkeytr{Zwiększanie} i \ealkeytr{zmniejszanie} bez kalkulatora}{\ealkey{Increasing} and \ealkey{decreasing} without a calculator}
    }
  \end{center}
  \vfill
\end{frame}
%================================================================

\qq[\LARGE]{\ealpara{Zwiększ $80$ o $10\%$}{\ealkey{Increase} $80$ by $10\%$}}
  {\steps{$10\%$ & $8$}\par\vspace{0.4em}$80+8=88$}
\qq[\LARGE]{\ealpara{Zmniejsz $60$ o $50\%$}{\ealkey{Decrease} $60$ by $50\%$}}
  {\steps{$50\%$ & $30$}\par\vspace{0.4em}$60-30=30$}
\qq[\LARGE]{\ealpara{Zwiększ $40$ o $25\%$}{\ealkey{Increase} $40$ by $25\%$}}
  {\steps{$25\%$ & $10$}\par\vspace{0.4em}$40+10=50$}
\qq[\LARGE]{\ealpara{Zmniejsz $90$ o $10\%$}{\ealkey{Decrease} $90$ by $10\%$}}
  {\steps{$10\%$ & $9$}\par\vspace{0.4em}$90-9=81$}
\qq[\LARGE]{\ealpara{Zwiększ \pounds $250$ o $20\%$}{\ealkey{Increase} \pounds $250$ by $20\%$}}
  {\steps{$10\%$ & \pounds $25$ \\ $20\%$ & \pounds $50$}\par\vspace{0.4em}
   \pounds $250+$ \pounds $50=$ \pounds $300$}
\qq[\LARGE]{\ealpara{Zmniejsz $45$ o $20\%$}{\ealkey{Decrease} $45$ by $20\%$}}
  {\steps{$10\%$ & $4.5$ \\ $20\%$ & $9$}\par\vspace{0.4em}$45-9=36$}
\qq[\LARGE]{\ealpara{Zwiększ $120$ o $5\%$}{\ealkey{Increase} $120$ by $5\%$}}
  {\steps{$10\%$ & $12$ \\ $5\%$ & $6$}\par\vspace{0.4em}$120+6=126$}
\qq[\LARGE]{\ealpara{Zmniejsz \pounds $64$ o $25\%$}{\ealkey{Decrease} \pounds $64$ by $25\%$}}
  {\steps{$25\%$ & \pounds $16$}\par\vspace{0.4em}
   \pounds $64-$ \pounds $16=$ \pounds $48$}
\qq[\LARGE]{\ealpara{Zwiększ $350$ o $30\%$}{\ealkey{Increase} $350$ by $30\%$}}
  {\steps{$10\%$ & $35$ \\ $30\%$ & $105$}\par\vspace{0.4em}$350+105=455$}
\qq[\LARGE]{\ealpara{Zmniejsz $480$ o $35\%$}{\ealkey{Decrease} $480$ by $35\%$}}
  {\steps{$10\%$ & $48$ \\ $30\%$ & $144$ \\ $5\%$ & $24$ \\ $35\%$ & $168$}%
   \par\vspace{0.4em}$480-168=312$}

%================================================================
\section{Finding the original value without a calculator}
\begin{frame}[plain]
  \vfill
  \begin{center}
    {\usebeamercolor[fg]{structure}\Huge\bfseries
    \ealpara{Znajdowanie \ealkeytr{wartości początkowej} bez kalkulatora}{Finding the \ealkey{original value} without a calculator}
    }
  \end{center}
  \vfill
\end{frame}
%================================================================

\qq[\Large][\large]{\ealpara{$10\%$ pewnej liczby to $7$.\\[0.3em]Jaka to liczba?}{$10\%$ of a number is $7$.\\[0.3em]What is the number?}}
  {\barmodel{10}{1}{$10\%$ is $7$}{$100\%$ is $?$}\par\vspace{0.9em}
   \steps{$10\%$ & $7$ \\ $100\%$ & $70$}}

\qq[\Large][\large]{\ealpara{$50\%$ pewnej liczby to $24$.\\[0.3em]Jaka to liczba?}{$50\%$ of a number is $24$.\\[0.3em]What is the number?}}
  {\barmodel{2}{1}{$50\%$ is $24$}{$100\%$ is $?$}\par\vspace{0.9em}
   \steps{$50\%$ & $24$ \\ $100\%$ & $48$}}

\qq[\Large][\large]{\ealpara{$20\%$ pewnej liczby to $12$.\\[0.3em]Jaka to liczba?}{$20\%$ of a number is $12$.\\[0.3em]What is the number?}}
  {\barmodel{10}{2}{$20\%$ is $12$}{$100\%$ is $?$}\par\vspace{0.9em}
   \steps{$20\%$ & $12$ \\ $10\%$ & $6$ \\ $100\%$ & $60$}}

\qq[\Large][\large]{\ealpara{$80\%$ pewnej liczby to $64$.\\[0.3em]Jaka to liczba?}{$80\%$ of a number is $64$.\\[0.3em]What is the number?}}
  {\barmodel{10}{8}{$80\%$ is $64$}{$100\%$ is $?$}\par\vspace{0.9em}
   \steps{$80\%$ & $64$ \\ $10\%$ & $8$ \\ $100\%$ & $80$}}

\qq[\Large][\large]{\ealpara{$30\%$ pewnej liczby to $18$.\\[0.3em]Jaka to liczba?}{$30\%$ of a number is $18$.\\[0.3em]What is the number?}}
  {\barmodel{10}{3}{$30\%$ is $18$}{$100\%$ is $?$}\par\vspace{0.9em}
   \steps{$30\%$ & $18$ \\ $10\%$ & $6$ \\ $100\%$ & $60$}}

\qq[\Large][\large]{\ealpara{$25\%$ pewnej liczby to $15$.\\[0.3em]Jaka to liczba?}{$25\%$ of a number is $15$.\\[0.3em]What is the number?}}
  {\barmodel{4}{1}{$25\%$ is $15$}{$100\%$ is $?$}\par\vspace{0.9em}
   \steps{$25\%$ & $15$ \\ $100\%$ & $60$}}

\qq[\Large][\large]{\ealpara{$15\%$ pewnej liczby to $12$.\\[0.3em]Jaka to liczba?}{$15\%$ of a number is $12$.\\[0.3em]What is the number?}}
  {\barmodel{20}{3}{$15\%$ is $12$}{$100\%$ is $?$}\par\vspace{0.9em}
   \steps{$15\%$ & $12$ \\ $5\%$ & $4$ \\ $100\%$ & $80$}}

\qq[\Large][\large]{\ealpara{Płaszcz jest \ealkeytr{obniżony} o $20\%$ w wyprzedaży.\\[0.3em]Teraz kosztuje \pounds $48$. Jaka była cena początkowa?}{A coat is \ealkey{reduced} by $20\%$ in a sale.\\[0.3em]It now costs \pounds $48$. What was the original price?}}
  {\barmodel{10}{8}{$80\%$ is \pounds $48$}{$100\%$ is $?$}\par\vspace{0.9em}
   \steps{$80\%$ & \pounds $48$ \\ $10\%$ & \pounds $6$ \\ $100\%$ & \pounds $60$}}

\qq[\Large][\large]{\ealpara{Cena wzrasta o $10\%$ do \pounds $88$.\\[0.3em]Jaka była cena przed wzrostem?}{A price rises by $10\%$ to \pounds $88$.\\[0.3em]What was the price before the rise?}}
  {\barmodel{11}{10}{$100\%$ is $?$}{$110\%$ is \pounds $88$}\par\vspace{0.9em}
   \steps{$110\%$ & \pounds $88$ \\ $10\%$ & \pounds $8$ \\ $100\%$ & \pounds $80$}}

\qq[\Large][\large]{\ealpara{Dom wzrasta na wartości o $30\%$ do \pounds $169\,000$.\\[0.3em]Ile był wart przed wzrostem?}{A house rises in value by $30\%$ to \pounds $169\,000$.\\[0.3em]What was it worth before the rise?}}
  {\barmodel{13}{10}{$100\%$ is $?$}{$130\%$ is \pounds $169\,000$}\par\vspace{0.9em}
   \steps{$130\%$ & \pounds $169\,000$ \\ $10\%$ & \pounds $13\,000$ \\ $100\%$ & \pounds $130\,000$}}

%================================================================
\section{Percentage word problems}
\begin{frame}[plain]
  \vfill
  \begin{center}
    {\usebeamercolor[fg]{structure}\Huge\bfseries
    \ealpara{Zadania tekstowe z \ealkeytr{procentami}}{\ealkey{Percentage} word problems}
    }
  \end{center}
  \vfill
\end{frame}
%================================================================

\qq[\Large][\Large]{\ealpara{Płaszcz kosztuje \pounds $40$.\\[0.3em]W wyprzedaży jest \ealkeytr{obniżony} o $10\%$.\\[0.3em]Jaka jest cena po obniżce?}{A coat costs \pounds $40$.\\[0.3em]In a sale it is \ealkey{reduced} by $10\%$.\\[0.3em]What is the sale price?}}
  {\steps{$10\%$ & \pounds $4$}\par\vspace{0.4em}
   \pounds $40-$ \pounds $4=$ \pounds $36$}

\qq[\Large][\Large]{\ealpara{$30\%$ klasy liczącej $30$ uczniów chodzi do szkoły pieszo.\\[0.3em]Ilu uczniów chodzi do szkoły pieszo?}{$30\%$ of a class of $30$ students walk to school.\\[0.3em]How many students walk to school?}}
  {\steps{$10\%$ & $3$ \\ $30\%$ & $9$}\par\vspace{0.4em}
   \ealpara{$9$ uczniów}{$9$ students}}

\qq[\Large][\Large]{\ealpara{Telefon kosztuje \pounds $250$ plus $20\%$ VAT.\\[0.3em]Jaka jest cena całkowita?}{A phone costs \pounds $250$ plus $20\%$ VAT.\\[0.3em]What is the total price?}}
  {\steps{$10\%$ & \pounds $25$ \\ $20\%$ & \pounds $50$}\par\vspace{0.4em}
   \pounds $250+$ \pounds $50=$ \pounds $300$}

\qq[\Large][\Large]{\ealpara{Sam zarabia \pounds $800$ miesięcznie i otrzymuje $5\%$ podwyżki.\\[0.3em]Jaka jest nowa miesięczna pensja Sama?}{Sam earns \pounds $800$ a month and is given a $5\%$ pay rise.\\[0.3em]What is Sam's new monthly pay?}}
  {\steps{$10\%$ & \pounds $80$ \\ $5\%$ & \pounds $40$}\par\vspace{0.4em}
   \pounds $800+$ \pounds $40=$ \pounds $840$}

\qq[\Large][\Large]{\ealpara{Kurtka kosztowała \pounds $60$, a teraz kosztuje \pounds $45$.\\[0.3em]Jaki jest \ealkeytr{procent} obniżki?}{A jacket cost \pounds $60$ and now costs \pounds $45$.\\[0.3em]What is the \ealkey{percentage} reduction?}}
  {\ealpara{Obniżka $=$ \pounds $15$}{Reduction $=$ \pounds $15$}\par\vspace{0.4em}
   $\dfrac{15}{60}=\dfrac{25}{100}=25\%$}

\qq[\Large][\Large]{\ealpara{Populacja miasta wzrasta z $8000$ do $8600$.\\[0.3em]Jaki jest \ealkeytr{procent} \ealkeytr{zwiększenia}?}{A town's population grows from $8000$ to $8600$.\\[0.3em]What is the \ealkey{percentage} \ealkey{increase}?}}
  {\ealpara{\ealkeytr{Zwiększenie} $=600$}{\ealkey{Increase} $=600$}\par\vspace{0.4em}
   $\dfrac{600}{8000}=\dfrac{7.5}{100}=7.5\%$}

\qq[\Large][\large]{\ealpara{Laptop kosztuje \pounds $480$ w wyprzedaży, po \ealkeytr{rabacie} $20\%$.\\[0.3em]Jaka była cena przed wyprzedażą?}{A laptop costs \pounds $480$ in a sale, after a $20\%$ \ealkey{discount}.\\[0.3em]What was the price before the sale?}}
  {\barmodel{10}{8}{$80\%$ is \pounds $480$}{$100\%$ is $?$}\par\vspace{0.9em}
   \steps{$80\%$ & \pounds $480$ \\ $10\%$ & \pounds $60$ \\ $100\%$ & \pounds $600$}}

\qq[\Large][\Large]{\ealpara{Maya zdobyła $34$ z $40$ w teście.\\[0.3em]Ben zdobył $82\%$. Kto wypadł lepiej?}{Maya scored $34$ out of $40$ in a test.\\[0.3em]Ben scored $82\%$. Who did better?}}
  {$\dfrac{34}{40}=\dfrac{85}{100}=85\%$\par\vspace{0.4em}
   \ealpara{$85\%>82\%$, więc Maya wypadła lepiej}{$85\%>82\%$, so Maya did better}}

\qq[\Large][\Large]{\ealpara{Telewizor kosztuje \pounds $1200$.\\[0.3em]Jest \ealkeytr{obniżony} o $15\%$, a następnie o kolejne $10\%$ nowej ceny.\\[0.3em]Jaka jest cena końcowa?}{A television costs \pounds $1200$.\\[0.3em]It is \ealkey{reduced} by $15\%$, then by a further $10\%$ of the new price.\\[0.3em]What is the final price?}}
  {\ealpara{
     \steps{$15\%$ z \pounds $1200$ & \pounds $180$ \\
             nowa cena & \pounds $1020$ \\
             $10\%$ z \pounds $1020$ & \pounds $102$}\par\vspace{0.4em}
      Cena końcowa $=$ \pounds $918$
   }{
     \steps{$15\%$ of \pounds $1200$ & \pounds $180$ \\
             new price & \pounds $1020$ \\
             $10\%$ of \pounds $1020$ & \pounds $102$}\par\vspace{0.4em}
      Final price $=$ \pounds $918$
   }}

\qq[\Large][\Large]{\ealpara{Samochód kosztuje \pounds $12\,000$.\\[0.3em]Traci $20\%$ swojej wartości w pierwszym roku i $15\%$ swojej wartości w drugim roku.\\[0.3em]Ile jest wart po dwóch latach?}{A car costs \pounds $12\,000$.\\[0.3em]It loses $20\%$ of its value in the first year and $15\%$ of its value in the second year.\\[0.3em]What is it worth after two years?}}
  {\ealpara{
     \steps{po roku $1$ & \pounds $12\,000\times 0.8=$ \pounds $9600$ \\
             po roku $2$ & \pounds $9600\times 0.85=$ \pounds $8160$}
   }{
     \steps{after year $1$ & \pounds $12\,000\times 0.8=$ \pounds $9600$ \\
             after year $2$ & \pounds $9600\times 0.85=$ \pounds $8160$}
   }}

%================================================================
\section{Harder percentages of a number}
\begin{frame}[plain]
  \vfill
  \begin{center}
    {\usebeamercolor[fg]{structure}\Huge\bfseries
    \ealpara{Trudniejsze \ealkeytr{procenty} liczby}{Harder \ealkey{percentages} of a number}
    }
  \end{center}
  \vfill
\end{frame}
%================================================================

\qq{\ealpara{$11\%$ z $200$}{$11\%$ of $200$}}
  {\steps{$10\%$ & $20$ \\ $1\%$ & $2$ \\ $11\%$ & $22$}}
\qq{\ealpara{$21\%$ z $400$}{$21\%$ of $400$}}
  {\steps{$10\%$ & $40$ \\ $20\%$ & $80$ \\ $1\%$ & $4$ \\ $21\%$ & $84$}}
\qq{\ealpara{$17.5\%$ z $40$}{$17.5\%$ of $40$}}
  {\steps{$10\%$ & $4$ \\ $5\%$ & $2$ \\ $2.5\%$ & $1$ \\ $17.5\%$ & $7$}}
\qq{\ealpara{$12.5\%$ z $80$}{$12.5\%$ of $80$}}
  {\steps{$10\%$ & $8$ \\ $2.5\%$ & $2$ \\ $12.5\%$ & $10$}}
\qq{\ealpara{$17.5\%$ z $240$}{$17.5\%$ of $240$}}
  {\steps{$10\%$ & $24$ \\ $5\%$ & $12$ \\ $2.5\%$ & $6$ \\ $17.5\%$ & $42$}}
\qq{\ealpara{$32\%$ z $150$}{$32\%$ of $150$}}
  {\steps{$10\%$ & $15$ \\ $30\%$ & $45$ \\ $1\%$ & $1.5$ \\ $2\%$ & $3$ \\ $32\%$ & $48$}}
\qq{\ealpara{$17.5\%$ z $68$}{$17.5\%$ of $68$}}
  {\steps{$10\%$ & $6.8$ \\ $5\%$ & $3.4$ \\ $2.5\%$ & $1.7$ \\ $17.5\%$ & $11.9$}}
\qq{\ealpara{$87.5\%$ z $96$}{$87.5\%$ of $96$}}
  {\steps{$100\%$ & $96$ \\ $10\%$ & $9.6$ \\ $2.5\%$ & $2.4$ \\ $12.5\%$ & $12$ \\ $87.5\%$ & $96-12=84$}}
\qq{\ealpara{$6.5\%$ z $400$}{$6.5\%$ of $400$}}
  {\steps{$1\%$ & $4$ \\ $6\%$ & $24$ \\ $0.5\%$ & $2$ \\ $6.5\%$ & $26$}}
\qq[\Huge][\LARGE]{\ealpara{$17.5\%$ z \pounds $312$}{$17.5\%$ of \pounds $312$}}
  {\steps{$10\%$ & \pounds $31.20$ \\ $5\%$ & \pounds $15.60$ \\
          $2.5\%$ & \pounds $7.80$ \\ $17.5\%$ & \pounds $54.60$}}

\end{document}